\section{The formal wedge obstruction}
\label{sec:wedge-obstruction}

We finish by comparing the canonical map with the map forced by the
proposed wedge. The argument uses naturality and chromatic acyclicity;
it places no compatibility requirement on a hypothetical equivalence
of spectra.

\begin{lemma}\label{lem:wedge-obstruction}
Suppose \(Z\simeq A\vee B\), \(L_{K(2)}B=0\), and
\(\pi_dL_0A=0\). Then \(\pi_d(\tau_Z)=0\).
\end{lemma}

\begin{proof}
Localization is exact and preserves finite sums, so
\(L_{K(2)}(A\vee B)\simeq L_{K(2)}A\).
Naturality with the summand inclusions identifies
\(\tau_{A\vee B}\) with \(\tau_A\) on \(L_0A\) and the zero map on
\(L_0B\). It vanishes in degree \(d\). Naturality with an arbitrary
equivalence \(A\vee B\simeq Z\) transports this conclusion to \(Z\).
\end{proof}

\begin{proof}[Proof of \cref{cor:strong}]
In \eqref{eq:proposed-wedge}, let \(A\) be the first line and \(B\)
the sum of the last two lines. Smashing for \(L_1\) and the
Bousfield-class description of \(L_1S\) give \(L_{K(2)}B=0\)
\cite[Theorem~5.1 and Proposition~5.3]{HS1999}.

The \(p\)-complete sphere is connective. Each term of the Moore tower
\(S_p^\wedge=\varprojlim_m S/p^m\) is connective, and the degree-zero
transition maps are surjective; the possible negative Milnor term
therefore vanishes. Rationalization preserves connectivity. Since
\(L_0L_2\simeq L_0\), the two summands of \(L_0A\) have no homotopy
in degree \(-3\). Thus \cref{lem:wedge-obstruction} gives
\(\pi_{-3}\tau_{\mathcal W_3}=0\).

Put \(T=L_{K(3)}S\) and \(X=L_2T\). The cofiber of the localization
map \(T\to X\) is \(E(2)\)-acyclic. The identity
\[
 \langle E(2)\rangle
       =\langle K(0)\vee K(1)\vee K(2)\rangle
\]
makes it both rationally and \(K(2)\)-acyclic. The naturality square
\[
\begin{tikzcd}[column sep=large,row sep=large]
 L_0T\arrow[r,"\tau_T"]\arrow[d,"\sim"']
      &L_0L_{K(2)}T\arrow[d,"\sim"]\\
 L_0X\arrow[r,"\tau_X"]&L_0L_{K(2)}X
\end{tikzcd}
\]
therefore has vertical equivalences. Any equivalence
\(X\simeq\mathcal W_3\) would force its lower horizontal map to vanish
on \(\pi_{-3}\), and hence force the same for \(\tau_T\).
This contradicts \cref{thm:main}. No step specifies the action of
the hypothetical equivalence on individual summands.
\end{proof}
